设 $\mathcal{C}$ 是范畴. 考虑以所有自函子 $T: \mathcal{C} \to \mathcal{C}$ 为对象, 以自函子之间的态射构成的函子范畴 $\mathcal{C}^{\mathcal{C}}$. 它自然地成为一个幺半范畴, 其乘法 $\otimes$ 是函子的合成运算, 幺元 $\munit$ 是恒等函子. 因为函子的合成满足严格结合律, 它还是严格幺半范畴.

\begin{definition}[单子和余单子]\label{def:monad}
	\index{danzi@单子 (monad)}
	\index{yudanzi@余单子 (comonad)}
	设 $\mathcal{C}$ 为范畴. 幺半范畴 $\mathcal{C}^{\mathcal{C}}$ 中的代数 (定义 \ref{def:algebra-monoidal}) 称为 $\mathcal{C}$ 上的单子. 对偶地, $\mathcal{C}^{\opp}$ 上的单子称为 $\mathcal{C}$ 上的余单子.
\end{definition}

对自函子范畴 $\mathcal{C}^{\mathcal{C}}$ 展开定义 \ref{def:algebra-monoidal}, 可知单子是资料 $(T, \mu, \eta)$, 其中 $T: \mathcal{C} \to \mathcal{C}$ 是函子, $\mu: T^2 \to T$ 和 $\eta: \identity_{\mathcal{C}} \to T$ 是态射, 使得下图交换
\begin{equation*}\begin{gathered}
	\begin{tikzcd}
		T \arrow[r, "{\eta T}"] \arrow[rd, "\identity"'] & T^2 \arrow[d, "\mu"] & T \arrow[l, "{T\eta}"'] \arrow[ld, "\identity"] \\
		& T &
	\end{tikzcd} \quad
	\begin{tikzcd}
		T^3 \arrow[r, "\mu T"] \arrow[d, "T\mu"'] & T^2 \arrow[d, "\mu"] \\
		T^2 \arrow[r, "\mu"'] & T.
	\end{tikzcd}
\end{gathered}\end{equation*}

相较于原定义, 此处不再需要结合约束, 而 $\eta \otimes \identity$ (或 $\identity \otimes \eta$) 则被翻译为 $\eta T$ (或 $T\eta$), 依此类推.

对偶地, 余单子 $(L, \delta, \epsilon)$ 由函子 $L: \mathcal{C} \to \mathcal{C}$, 态射 $\delta: L \to L^2$ 和 $\epsilon: L \to \identity_{\mathcal{C}}$ 组成, 条件是使下图交换.
\begin{equation*}\begin{gathered}
	\begin{tikzcd}
		L  & L^2 \arrow[l, "{\epsilon L}"'] \arrow[r, "{L\epsilon}"] & L \\
		& L \arrow[lu, "\identity"] \arrow[ru, "\identity"'] \arrow[u, "\delta"] &
	\end{tikzcd} \quad
	\begin{tikzcd}
		L^3 & L^2 \arrow[l, "\delta L"'] \\
		L^2 \arrow[u, "L\delta"] & L \arrow[u, "\delta"'] \arrow[l, "\delta"]
	\end{tikzcd}
\end{gathered}\end{equation*}

我们习惯将资料 $(T, \mu, \eta)$ (或 $(L, \delta, \epsilon)$) 简记为 $T$ (或 $L$). 单子和余单子的主要来源是伴随对.

\begin{example}[伴随对确定单子]\label{eg:adjunction-monad}
	考虑一对伴随函子
	\[\begin{tikzcd}
		F: \mathcal{C} \arrow[shift left, r] & \mathcal{D}: G \arrow[shift left, l],
	\end{tikzcd}\]
	以及相应的单位 $\eta: \identity_{\mathcal{C}} \to GF$ 和余单位 $\varepsilon: FG \to \identity_{\mathcal{D}}$ 态射. 今将定义 $\mathcal{C}$ 上的单子 $(T, \mu, \eta)$ 和 $\mathcal{D}$ 上的余单子 $(L, \delta, \varepsilon)$ 如下.
	\[\begin{array}{|c|c|}\hline
		T := GF & L := FG \\
		\mu := \left[ GFGF \xrightarrow{G\varepsilon F} GF \right] & \delta := \left[ FG \xrightarrow{F\eta G} FGFG \right] \\
		\eta : \identity_{\mathcal{C}} \to GF & \varepsilon: FG \to \identity_{\mathcal{D}} \\ \hline 
	\end{array}\]

	基于对偶性, 就 $(T, \mu, \eta)$ 的情形验证单子的公理即可. 首先是关于幺元的图表
	\[\begin{tikzcd}[column sep=large]
		GF \arrow[r, "{\eta GF}"] \arrow[equal, rd] & GFGF \arrow[d, "{G\varepsilon F}"] & GF \arrow[l, "{GF\eta}"'] \arrow[equal, ld] \\
		& GF &
	\end{tikzcd}\]
	其交换性来自三角等式 $(G\varepsilon)(\eta G) = \identity_G$ 和 $(\varepsilon F)(F\eta) = \identity_F$. 结合律图表是
	\[\begin{tikzcd}[column sep=large]
		GFGFGF \arrow[d, "{G\varepsilon FGF}"'] \arrow[r, "{GFG\varepsilon F}"] & GFGF \arrow[d, "G\varepsilon F"] \\
		GFGF \arrow[r, "G\varepsilon F"'] & GF
	\end{tikzcd}\]
	它是交换的: 两路合成同样图解为
	\[\begin{tikzcd}
		\mathcal{C} \arrow[r, "F"] & \mathcal{D} \arrow[r, "G"] \arrow[rr, bend right=60, ""{name=A}, "{\identity}"'] & \mathcal{C} \arrow[r, "F"] \arrow[Rightarrow, to=A, "\varepsilon"] &
		\mathcal{D} \arrow[r, "G"] \arrow[rr, bend right=60, ""{name=B}, "{\identity}"'] & \mathcal{C} \arrow[r, "F"] \arrow[Rightarrow, to=B, "\varepsilon"] & \mathcal{D} \arrow[r, "G"] & \mathcal{C}
	\end{tikzcd}\]
	只是两个弓形区域的合成次序不同, 其产物则相同; 这是态射纵横合成的互换律 \cite[引理 2.2.7]{Li1} 的一则特例.
\end{example}

由于自函子作用在 $\mathcal{C}$ 上, 我们希望在 $\mathcal{C}$ 中探讨在单子 $T$ 作用下的``模''. 由于历史的原因, 这种结构也被称为 $T$-代数.

\begin{definition}[S.\ Eilenberg, J.\ C.\ Moore]\label{def:T-Alg}
	\index{mo}
	设 $(T, \mu, \eta)$ 是 $\mathcal{C}$ 上的单子. 所谓 $T$-模, 系指资料 $(M, a)$, 其中 $M \in \Obj(\mathcal{C})$ 而 $a$ 是 $\mathcal{C}$ 的态射 $T(M) \to M$, 使得下图交换
	\[\begin{tikzcd}
		T^2 (M) \arrow[r, "Ta"] \arrow[d, "\mu_M"'] & T(M) \arrow[d, "a"] \\
		T(M) \arrow[r, "a"'] & M
	\end{tikzcd} \quad \begin{tikzcd}
		M \arrow[r, "\eta_M"] \arrow[rd, "{\identity_M}"'] & T(M) \arrow[d, "a"] \\
		& M .
	\end{tikzcd}\]

	所有 $T$-模构成范畴 $\mathcal{C}^T$: 从 $(M, a)$ 到 $(M', a')$ 的态射定为 $\mathcal{C}$ 中使得下图交换的态射 $f: M \to M'$
	\[\begin{tikzcd}
		T(M) \arrow[r, "a"] \arrow[d, "Tf"'] & M \arrow[d, "f"] \\
		T(M')  \arrow[r, "{a'}"'] & M' .
	\end{tikzcd}\]

	对偶地, 设 $(L, \delta, \epsilon)$ 是 $\mathcal{C}$ 上的余单子, 所谓 $L$-余模是资料 $(N, b)$, 其中 $b$ 是态射 $N \to L(N)$, 条件是有与先前相对偶的交换图表. 我们将 $L$-余模范畴记为 $\mathcal{C}^L$.
\end{definition}

关于 $T$-模的两个交换图表应当分别被设想为 $T$ 作用下的结合律和幺元律. 今后的例子和性质主要针对单子及其上的模加以陈述, 余单子和余模的版本纯然是对偶的.
